\documentclass[11pt]{article}
\usepackage[margin=1in]{geometry}
\usepackage{lmodern,microtype,amsmath,amssymb,amsthm,booktabs,tabularx,enumitem,needspace}
\usepackage[colorlinks=true,linkcolor=blue,urlcolor=blue]{hyperref}
\interfootnotelinepenalty=10000
\setlength{\parindent}{0pt}
\setlength{\parskip}{5pt}
\setlist{nosep,leftmargin=*}
\newtheorem{proposition}{Proposition}
\title{Housing allocation and constrained efficiency}
\author{Discussion note}
\date{September 5, 2026}
\begin{document}
\maketitle
\vspace{-1em}
The two planner comparisons answer different questions. A direct-allocation
planner can redistribute housing and consumption subject to real feasibility.
A constrained planner redistributes cash, after which households choose and
markets clear. In both comparisons, the housing technology and physical tenure
segmentation are retained. Fertility and the population path are held fixed
for the housing-welfare question.

\textbf{Assessment.} A simple local Pareto improvement toward young owners is
proved for the direct-allocation benchmark under the value comparison below.
Constrained inefficiency of the complete OLG economy remains unproved. A
conditional market exercise supplies both an improving cash-transfer example
and an obstruction, so the stronger conclusion does not follow from the
financing gap alone. The positive example gives more housing to old households.
Its future markets, tenure changes, and rental-intermediary owners are not
covered. It should remain an internal check.

\section{Direct allocation}

\textbf{Households and feasible changes.} Consider a young owner $i$ and an old
owner $j$. The young household has adult consumption $x_i=c_i-\chi n_i>0$ and
adult space $s_i=h_i-\kappa n_i>0$. Its current utility is
$\log x_i+\alpha\log s_i+\vartheta_t\log n_i$. The old household has utility
$\log c_j^2+\gamma\log h_j^2+\omega_B\log e_j$. Keep each household's tenure,
fertility, estate, and every future consumption and housing allocation fixed.
The tenure-taste term and the young household's continuation utility therefore
remain unchanged.

The housing stock is fixed. Rental and owner size bounds remain in place;
there are no newly assumed separate rental and owner stocks. The proposed
transfer is owner to owner: the young owner's housing rises by $\epsilon$ and
the old owner's occupancy falls by the same amount. The young owner must have
room below $h_O^{\max}$, and the old owner must have positive housing to give up.
A real cost $C(\epsilon)$ uses current goods, with $C(0)=0$ and finite right
derivative $C'(0+)=K\geq0$.

The value of a housing unit in each household's own consumption is:
\begin{equation}
 MV_i^Y=\frac{\alpha x_i}{s_i},\qquad
 MV_j^O=\frac{\gamma c_j^2}{h_j^2}.
\end{equation}
These ratios compare willingness to exchange consumption for housing. Unequal
raw housing marginal utilities would not, by themselves, establish inefficiency.

\Needspace{9\baselineskip}
\begin{proposition}[An improving housing reallocation]
Suppose the owner-to-owner transfer is physically feasible and the planner may
redistribute current goods and settle financial claims while relaxing private
borrowing limits. If $MV_i^Y>MV_j^O+K$, sufficiently small positive transfers of
housing from the old owner to the young owner are Pareto improvements. The old
owner is indifferent, the young owner gains, and every other household and
outside creditor can retain its allocation or contractual payment.
\end{proposition}

\begin{proof}
Exactly compensating the old owner requires extra current consumption:
\begin{equation}
 D_j(\epsilon)=c_j^2\left[
 \left(\frac{h_j^2}{h_j^2-\epsilon}\right)^\gamma-1\right].
\end{equation}
Give the old owner $D_j(\epsilon)$ and reduce the young owner's current
consumption by $D_j(\epsilon)+C(\epsilon)$. Housing and current goods balance,
and the old owner's utility is unchanged. The young owner's lifetime gain is:
\begin{equation}
 \Delta U_i=\log\left(1-\frac{D_j(\epsilon)+C(\epsilon)}{x_i}\right)
 +\alpha\log\left(1+\frac{\epsilon}{s_i}\right).
\end{equation}
All future utility arguments are unchanged. Since $D_j'(0)=MV_j^O$, the right
derivative is $(MV_i^Y-MV_j^O-K)/x_i>0$. Sufficiently small transfers preserve
positive consumption and respect the physical bounds.
\end{proof}

\textbf{Connection to the competitive equilibrium.} An old owner with slack
retention and estate-composition constraints satisfies $MV_j^O=q r_t$. A young
owner with positive saving, slack future retention and estate bounds, and a
strictly positive financing wedge satisfies $MV_i^Y>q r_t$, provided its physical
size cap is slack. Such a pair therefore gives the desired inefficiency result
at zero real cost, or whenever cost is smaller than the gap. The existence of
eligible pairs is a condition, not a consequence of age alone. A positive
aggregate reallocation requires a positive measure of such pairs.

\textbf{What is established.} There exists an improving allocation with more
housing for the eligible young households. This does not select the whole
economy's first-best allocation or imply that every choice of welfare weights
raises aggregate young housing. No welfare weights are needed for the local
Pareto argument.

\section{Resources, titles, and the first-best benchmark}

The planner's aggregate resource constraint counts goods, the physical housing
stock, real transfer costs, and claims on the rest of the world. A housing sale
transfers goods between counterparties; it does not destroy goods equal to the
house price. Claims held outside the economy belong in the external account,
and property taxes and rebates are consolidated. The world bond price $q$
remains relevant to trade across dates. Private borrowing limits
may be relaxed; outside creditors' payments and external solvency must still
be respected.

The local construction avoids a new aggregate financing assumption: current
goods balance, every estate expenditure remains fixed, and the net external
asset path is unchanged. It also admits an exact check against the original
household accounts. Write $u_t=q r_t=(1+q\tau^p)P_t-qP_{t+1}$ and
$L=D_j(\epsilon)-u_t\epsilon$. Change the young owner's saving account by:
\begin{equation}
 \Delta a_i'=\phi_tP_t\epsilon-qP_{t+1}\epsilon.
\end{equation}
\Needspace{5\baselineskip}
The change in its next-date resources, including title liquidation, is then:
\begin{equation}
 \Delta(a_{i,t+1}+P_{t+1}H_i)
 =\frac{\Delta a_i'}q-\frac{\phi_tP_t\epsilon}q+P_{t+1}\epsilon=0.
\end{equation}
The original young and old budgets hold with net current transfers $-L-C$ to
the young owner and $L$ to the old owner, and $C$ used by the transfer activity.
The seller replaces the estate's missing title proceeds with a bond. Its bond
purchase exactly offsets the buyer's net financial adjustment. Property-tax
reserves also offset because total taxable housing is unchanged.

The young owner sells the additional title next period, when the old owner's
estate would have sold it. Later housing availability is therefore unchanged.
The seller's lower occupancy relaxes its original retention and estate bounds.
The buyer's next-period retention bound relaxes, while its old-age occupancy
and estate remain fixed. If the saving-account adjustment would make private
saving negative, the planner must provide the permitted alternative finance;
the construction does not claim that all private financial restrictions hold.

\textbf{The remaining first-best choice.} A complete planner program needs a
resource interpretation of the estate condition $e\geq P_{t+1}h^2$. The retained
house is physical, but its lower bound in goods uses a market valuation. We
must specify how title proceeds and goods estates are settled in the global
planner problem. The origin of entrant wealth and initial external claims
must also be accounted for once. The local proof preserves these objects and
verifies settlement at the reference price path, so it does not require
choosing a new planner valuation rule.

\textbf{The allocation figure.} A valid crossing-value diagram can be drawn
along the old household's indifference curve. At transfer $\epsilon$, plot:
\begin{align}
 MV_i^Y(\epsilon)&=
 \frac{\alpha[x_i-D_j(\epsilon)-C(\epsilon)]}{s_i+\epsilon},\\
 MV_j^O(\epsilon)&=
 \frac{\gamma[c_j^2+D_j(\epsilon)]}{h_j^2-\epsilon}.
\end{align}
With a smooth increasing convex cost, the first falls and the second plus
$C'(\epsilon)$ rises. Their crossing, if reached before a physical bound, is
the best allocation for this pair holding the old owner's utility and all
other allocations fixed. It is not the entire economy's first best. The area
between the curves is not an exact utility loss; the utility derivative is
the value gap divided by the young owner's remaining adult consumption.

\section{Cash transfers followed by markets}

\textbf{The proposed instrument.} Let $\ell_i$ be a net cash transfer paid before
housing purchases; negative transfers are cash taxes. Balance the transfers
across the specified recipients and donors. Retain the financed share $\phi_t$
and all physical housing restrictions. A young owner's purchase constraint is
then $(1-\phi_t)P_th_i\leq b_i+\ell_i$, and the same $\ell_i$ enters its goods
budget once. Ordinary property-tax rebates still arrive after closing and
have their original separate budget. This is a proposed transfer permission;
it does not include public loans or compulsory future repayment.

The constrained planner must anticipate the resulting private choices and
market prices. This is the distinction in Bisin's constrained-efficiency
formulation. It does not imply that a borrowing distortion alone guarantees
an improving transfer scheme.\footnote{Alberto Bisin, \emph{General equilibrium
theory}, November 4, 2014, Section 3.2.4.
\href{https://bpb-us-e1.wpmucdn.com/wp.nyu.edu/dist/c/16384/files/2019/11/Lecture-notes-Oct2014.pdf}{NYU lecture notes}.}

\Needspace{7\baselineskip}
\textbf{A useful conditional test.} First fix future prices and rebates, the
inherited old states, fertility, and the local tenure and constraint branches.
At a candidate current price $P$, calculate the exact cash transfer
$\ell_i^c(P)$ needed to leave each household at its baseline utility. Let $G(P)$
be the sum of those transfers, weighted by household masses, and let $Q^c(P)$
be their total housing demand after compensation. At the baseline price $P_0$,
$G(P_0)=0$ and $Q^c(P_0)=\bar H$.

On the branches examined, utility increases with cash and housing demand is
nondecreasing in cash. A balanced Pareto reform must therefore satisfy:
\begin{equation}
 G(P)\leq0,\qquad Q^c(P)\leq\bar H.
\end{equation}
The first inequality says compensation is affordable. The second says that
compensating everybody must not already exhaust more than the housing stock;
additional welfare gains cannot reduce housing demand on these branches.

\begin{proposition}[A local obstruction for cash transfers]
In this conditional comparison, suppose exact compensation and demands are
smooth on the maintained branches. If $G'(P_0)>0$ and
$Q^{c\,\prime}(P_0)<0$, no sufficiently close balanced cash-transfer equilibrium
makes all the compared households weakly better off and any strictly better off.
\end{proposition}

\begin{proof}
A higher price requires positive net compensation, violating the balanced
transfer budget. A lower price raises compensated housing demand above the
stock, violating housing clearing. At the unchanged price, every household
requires a nonnegative transfer; budget balance forces all transfers to zero.
Strict derivative signs preserve these arguments in a neighborhood.
\end{proof}

The obstruction can occur with a positive property tax, an interior old owner,
and a strictly constrained young owner whose housing value exceeds the old
owner's. An independently checked construction gives
$G'(P_0)=0.60513>0$ and $Q^{c\,\prime}(P_0)=-0.09876<0$.
It satisfies the original current-date household equations and housing and tax
accounts. Its inherited old state is specified directly; it is not a claimed
stationary OLG equilibrium. The example disproves an automatic inference from
the financing gap to an improving cash transfer in this conditional comparison.
It does not prove constrained efficiency of the full OLG economy.

\textbf{A conditional improvement with the opposite housing direction.} In the
saved illustrative allocation, the same calculation instead gives
$G'(P_0)=0.52065>0$ and $Q^{c\,\prime}(P_0)=0.17244>0$. A small price reduction
releases enough resources after compensation to improve the two old household
groups. The exact construction keeps both young groups at their baseline
lifetime utility, lets old-owner demand clear housing, and gives the remaining
cash to old renters. It taxes the young groups and transfers cash to the old.
The private housing choices move toward old households.

The direction follows from the old household's compensated demand. When its
retention and estate bounds are slack, its housing rises as the current user
cost falls along an indifference curve; a utility gain raises its housing
further. Since affordable compensation requires a lower price in this case,
old housing rises. With renters at their physical caps and fixed total housing,
young-owner housing must fall. Thus an improving transfer among these household
groups does not implement the direct planner's old-to-young reallocation.

These calculations are internal checks. The positive construction covers four
current household groups, with their lifetime continuation values evaluated
at fixed future prices. It is not a complete-economy Pareto result.

\section{What is needed for the full constrained comparison}

\textbf{Future markets.} A current transfer changes young saving, title holdings,
and subsequent housing and estates even when young lifetime utility is held
constant. The next old distribution consequently changes. Future prices and
property-tax rebates must clear that changed economy, and all later cohorts
must be included in the welfare test. Fixing fertility preserves their masses;
it does not preserve their prices or welfare.

\textbf{Initial ownership.} Rental intermediaries earn their prescribed return
on units bought at the new price. That does not compensate the owners of titles
held before an unexpected price change. The household equations do not identify
the ultimate recipients of all initial intermediary and estate claims. A price
fall can reduce those claims' value. We must state who owns them and whether
their welfare belongs in the comparison before claiming a Pareto improvement.
Zero profits on new purchases do not settle that question.

\textbf{Transfer powers.} The initial cash-transfer class is narrower than an
arbitrary sequence of household-specific compensation payments. We must choose
whether the planner has one transfer date or a committed sequence of separately
balanced transfers, and which characteristics it may use for targeting. Adding
public bridge credit or future compulsory repayment is a different instrument.
Any positive theorem must name that permission explicitly.

The next research question is therefore specific: under a stated ownership
account and transfer schedule, can a fully clearing OLG path make every covered
household weakly better off? The direct-allocation proof supplies the housing
mechanism. It does not answer this stronger question. The conditional examples
show why both the fiscal permission and equilibrium response matter.

\section{Verification and working record}

The direct result was checked by deriving exact compensation from the original
preferences, reconciling the dated accounts and physical restrictions, and
independently testing finite transfers and boundary failures. The independent
check covers 27 constructed household perturbations with positive and zero
property taxes. It includes 18 gains when cost is below the gap and nine losses
when cost exceeds it. The maximum accounting error is $1.78\times10^{-15}$.
These are household checks, not equilibrium or calibration estimates.

The cash-transfer argument was checked by deriving compensation from the
original budgets, proving the finite-neighborhood obstruction, and constructing
exact finite transfers in the saved conditional illustration. Independent lead
calculations reproduce the positive example at price reductions of $10^{-4}$,
$10^{-5}$, and $10^{-6}$, with maximum residual $3.92\times10^{-16}$, and verify
the obstruction's household budgets, first-order conditions, and strict bounds.
The examples do not verify future market clearing or the welfare of omitted
initial title owners.

Two Astra agents at maximum reasoning wrote separate bounded reviews. The lead
checked their arguments and independently reproduced the numerical algebra.
The reports and checks are indexed in
\path{output/model/simplified_olg_amendments/README.md}.
Reproduce the lead checks from the repository root with:
\begin{quote}\small
\texttt{python3}\ \path{code/model/tools/verify_simplified_olg_planner_benchmarks.py}
\end{quote}
No author manuscript, calibration specification, or model solver was changed.

\end{document}
